\section{Fine-Grained Visual Classification}

Fine-grained visual classification (FGVC) refers to the task of distinguishing subordinate categories that belong to the same general visual category, such as bird species, aircraft variants, or car models \cite{Fine-grainedSurvey}. Compared with generic object recognition, FGVC is more challenging because the visual differences between classes are often subtle and localized, while images from the same class may still vary significantly due to pose, viewpoint, illumination, background, and object deformation \cite{Fine-grainedSurvey, Zhou_2016_CVPR, TransIFC}. In bird species recognition, two species may share nearly identical global shapes and color distributions, differing only in minute visual traits such as bill shape, wing pattern, crown color, or breast pattern \cite{wah2011cub,HierarchicalFeatureAtt, He_2017}. Consequently, effective FGVC systems must be capable of learning representations that simultaneously capture global object structure and localized discriminative details \cite{HierarchicalFeatureAtt}.

Early FGVC studies often relied on part annotations, bounding boxes, or explicitly designed part localization mechanisms to compare corresponding regions across images. For example, part-based R-CNN methods learn object and part detectors and use geometric constraints to construct pose-normalized representations for fine-grained categorization \cite{zhang2014part}. More recent deep learning methods reduce the dependency on manual part supervision by learning discriminative regions or feature interactions directly from image-level labels. Representative approaches include bilinear feature modeling \cite{lin2015bilinear}, recurrent attention \cite{fu2017racnn}, and region proposal-based attention mechanisms \cite{yang2018learning}. These methods demonstrate that discovering local discriminative evidence is central to fine-grained recognition. However, many conventional FGVC methods are formulated as single-level fine-grained classification problems, where the output is only the finest category and higher-level taxonomic consistency is not explicitly modeled \cite{Fine-grainedSurvey, park2025visually}.

This thesis focuses on fine-grained bird recognition in a hierarchical setting. Instead of predicting only the species label, the model predicts labels at multiple taxonomic levels, namely order, family, and species. This formulation is particularly relevant for bird classification, where the taxonomy provides useful semantic structure and where mistakes between closely related species are common. The following subsections review the CUB-200-2011 dataset, discriminative region learning in FGVC, and diversity regularization methods that motivate the attribute-guided learning strategy used in this thesis.


\subsection{Bird Species Recognition and the CUB-200-2011 Dataset}

The Caltech-UCSD Birds-200-2011 dataset \cite{wah2011cub}, commonly known as CUB-200-2011, is one of the most widely used benchmarks for fine-grained bird image classification. It contains 200 bird species and 11,788 images, with annotations including bounding boxes, part locations, and 312 binary visual attributes. These attributes describe semantic visual properties of birds, such as color, shape, and pattern, and are organized around bird parts such as the bill, wing, breast, belly, tail, head, and leg.

CUB-200-2011 is particularly suitable for studying attribute-guided fine-grained classification for several reasons. First, the species are visually similar and often differ only in localized traits. Second, the dataset provides rich attribute annotations that can be used as auxiliary semantic supervision. Third, the bird taxonomy can be organized into a hierarchy of order, family, and species, making it appropriate for hierarchical image classification. In this thesis, the CUB-200-2011 hierarchy follows the H-CAST setting with 13 orders, 38 families, and 200 species \cite{park2025visually}.

The attributes in CUB are useful because they provide human-interpretable descriptions of subtle visual cues. For instance, two bird species may have a similar body shape and overall color, but differ in bill shape, wing color, crown pattern, or belly pattern. These attributes can guide a model toward fine-grained traits that may be difficult to learn from class labels alone, consistent with prior attribute-guided and attribute-aware attention methods for fine-grained recognition \cite{liu2017localizing, han2018attribute}. In this thesis, CUB attributes are not used as manual inputs during inference. Instead, they are used during training as auxiliary supervision for learning attribute-oriented representations. This design preserves the standard image-only inference setting while exploiting additional semantic information available in the training data.

Nevertheless, using all available attributes may not always be optimal. Some attributes may be redundant, weakly discriminative, or irrelevant for distinguishing a particular group of sibling species. This motivates the study of attribute selection strategies, including the proposed Hard-Sibling Discriminative Attribute Selection (HSDAS), which aims to identify attributes that are useful for distinguishing taxonomically related and visually similar bird classes.


\subsection{Discriminative Region Learning in FGVC}

A central challenge in FGVC is learning where to look. Since fine-grained categories are often separated by small local differences, many methods aim to identify discriminative object regions or part-level patterns. Bilinear CNN models, for example, model pairwise interactions between local convolutional features and have been shown to be effective for texture-rich and fine-grained recognition tasks \cite{lin2015bilinear}. Although bilinear pooling does not explicitly require part annotations, it can capture local feature interactions that are useful for distinguishing visually similar categories.

Another line of work uses attention mechanisms to locate informative regions. RA-CNN recursively learns region attention and region-based feature representations at multiple scales, allowing the model to progressively focus on discriminative image regions \cite{fu2017racnn}. Similarly, NTS-Net introduces a Navigator-Teacher-Scrutinizer framework, where informative regions are proposed and then scrutinized for final fine-grained prediction \cite{yang2018learning}. These methods highlight the importance of localized evidence for fine-grained recognition.

Despite their effectiveness, many discriminative region learning methods are designed for flat species-level classification. They aim to improve the prediction of the finest label, but they do not explicitly model taxonomy-level consistency. In contrast, hierarchical classification requires the model to produce predictions at multiple semantic levels. For example, a species prediction should be compatible with its corresponding family and order. This thesis therefore builds on the visually consistent hierarchical classification framework of H-CAST, which uses segment-based visual grouping to align fine-to-coarse predictions \cite{park2025visually}.

The proposed approach differs from conventional discriminative region learning in two ways. First, it uses hierarchical visual representations from H-CAST rather than a flat image classifier. Second, it incorporates bird attributes as auxiliary supervision, encouraging the model to learn representations related to semantic visual traits. Instead of directly predicting attributes as the final task, the model uses attribute-oriented features to support species-level classification while preserving hierarchical predictions at the family and order levels.


\subsection{Loss-based Diversity Regularization}

In addition to architecture-based region learning, another direction in FGVC encourages discriminative and diverse feature learning through loss functions. Mutual-Channel Loss (MC-Loss) is a representative example of this direction \cite{Chang_2020}. MC-Loss was proposed to learn localized discriminative features without requiring explicit part or bounding-box annotations. It consists of two main components: a discriminality component and a diversity component. The discriminality component encourages feature channels associated with a class to be individually discriminative, while the diversity component encourages those channels to focus on different spatial regions.

The intuition behind the diversity component is that a fine-grained object cannot usually be recognized from a single local cue. For bird recognition, different parts such as the head, wing, breast, bill, and tail may all provide complementary evidence. If all feature channels attend to the same most salient region, the model may ignore other useful visual traits. By encouraging feature channels to be spatially diverse, MC-Loss helps the network discover multiple discriminative regions and reduce redundant feature responses \cite{Chang_2020}.

This idea is relevant to attribute-guided bird classification. In the proposed model, different attribute groups may correspond to different semantic parts or visual properties. For example, bill shape, crown color, wing pattern, and tail shape should ideally provide complementary information rather than collapsing onto the same image region. Therefore, this thesis adopts the diversity principle from MC-Loss and adapts it to the attribute-guided H-CAST setting.

It is important to note that the diversity regularization used in this thesis is not identical to the original MC-Loss. MC-Loss applies diversity regularization to class-specific convolutional feature channels. In contrast, this thesis applies a diversity-inspired regularization term to attribute-group attention patterns. The goal is to encourage different attribute groups or macro-parts to attend to complementary visual regions. Thus, the proposed loss should be understood as an adaptation of the MC-Loss diversity idea, rather than a direct reuse of the original MC-Loss formulation.

By combining attribute supervision with diversity regularization, the model is encouraged to learn attribute-oriented representations that are both discriminative and complementary. This is especially important for CUB-200-2011, where species-level differences often depend on multiple subtle visual attributes. The effect of this diversity regularization is evaluated in later experimental chapters through ablation studies with different loss weights.